\section{Requirements and Resources}
\label{sec:requirements}

This section lists the resources the project needs and the key functional
requirements it must meet. The stack uses familiar, well-documented tools.

\subsection{Software}

\begin{itemize}[leftmargin=1.4em,itemsep=0.2em]
\item \textbf{Backend:} a web application framework with a relational database.
\item \textbf{Vulnerable service:} a separate, isolated web application, with
      OWASP Juice Shop as a reference for challenge design~\cite{Owasp_juiceshop}.
      Challenge scope follows the OWASP Top~10~\cite{Owasp2021top10}.
\item \textbf{Scoring and gamification:} CTFd as the reference model for scoring,
      hints, and leaderboards~\cite{Ctfd_docs}.
\item \textbf{Machine learning:} a standard machine-learning library for the simple,
      interpretable at-risk model and next-challenge recommender.
\item \textbf{Deployment and tooling:} version control for the source, and
      \LaTeX{} for the report.
\end{itemize}

\subsection{Hardware}

\begin{itemize}[leftmargin=1.4em,itemsep=0.2em]
\item Development laptops for the team.
\item A server or virtual machine on the campus network to host the platform and the
      isolated vulnerable service (roughly 4 vCPUs, 8~GB RAM, 40~GB disk for a class
      pilot).
\end{itemize}

\subsection{Data sources}

The platform generates its own data: every attempt (payload, solved, accuracy, time,
hint depth, points) is logged and is the data the ML model learns from. No external
personal dataset is required. Challenge content is authored by the team using the
OWASP Top~10~\cite{Owasp2021top10} and OWASP Juice Shop~\cite{Owasp_juiceshop} as
references. Pilot evaluation data (pre/post skill scores and an engagement survey) is
collected from a consenting student cohort.

\subsection{Lab and deployment facilities}

The vulnerable service is network-isolated to the campus network and shares neither
credentials nor a data store with the platform, so exploits cannot reach real
records. A test environment is used for development, and a
supervised lab session hosts the evaluation pilot.

\subsection{Key functional requirements}

\begin{itemize}[leftmargin=1.4em,itemsep=0.2em]
\item Host a vulnerable app covering OWASP Top~10 families, each verified by a
      server-side state change (not a guessable flag). (O1)
\item Score every attempt automatically from completion, accuracy, and time. (O3)
\item Award points, badges, levels, and achievements, with a peer-banded
      leaderboard. (O2)
\item Provide a hint ladder and learning resources that unlock on struggle. (O4)
\item Train a simple ML model to predict at-risk learners and recommend the next
      challenge. (O4, O5)
\item Present an instructor dashboard of per-learner and cohort progress, completion,
      and scores. (O5)
\end{itemize}
